\section{Эксперименты}

\subsection{Методология и план экспериментов}

Для валидации предложенного подхода к неполной верификации 
с Tensor Parallelism были поставлены три группы экспериментов, 
каждая из которых проверяет ключевое свойство реализации:

\begin{enumerate}
    \item \textbf{Эксперимент на потребление памяти} (OOM-тест).
    Цель~--- продемонстрировать, что TP позволяет верифицировать 
    модели, не помещающиеся в память одного GPU. Сравнивается 
    пиковое потребление GPU-памяти при CROWN-верификации 
    на одном GPU и при распределении на два GPU.
    
    \item \textbf{Эксперимент на корректность bounds}.
    Цель~--- убедиться, что TP-реализация вычисляет 
    \textit{звуковые} (sound) bounds: нижняя граница $lb_{\text{tp}}$ 
    не превышает эталонную $lb_{\text{ref}}$, 
    а верхняя $ub_{\text{tp}}$ не меньше эталонной $ub_{\text{ref}}$.
    Сравниваются результаты TP с эталонным однопроцессорным CROWN 
    на моделях различной глубины.
    
    \item \textbf{Эксперимент на реальных моделях} (VNN-COMP).
    Цель~--- проверить работоспособность автоматического 
    шардирования на моделях в формате ONNX из бенчмарка 
    VNN-COMP 2022\cite{arXiv2212VNNCOMP2022}, подтвердив 
    применимость подхода за пределами вручную написанных 
    PyTorch-моделей.
\end{enumerate}

\textbf{Аппаратная конфигурация.}
Все эксперименты проводились на сервере со следующей конфигурацией:
\begin{itemize}
    \item 2$\times$ NVIDIA A40 (48~ГБ VRAM каждый, архитектура Ampere);
    \item межGPU-связь через NVLink;
    \item бэкенд коллективных операций: NCCL;
    \item фреймворк: PyTorch 2.x, запуск через \texttt{torchrun --nproc\_per\_node=2}.
\end{itemize}

\subsection{Эксперимент 1: потребление GPU-памяти}

\subsubsection{Архитектура модели}

В качестве объекта верификации использовалась полносвязная сеть 
с одним скрытым слоем:
$$f(\mathbf{x}) = W_2 \cdot \mathrm{ReLU}(W_1 \mathbf{x} + \mathbf{b}_1) + \mathbf{b}_2,$$
где $W_1 \in \mathbb{R}^{H \times D}$, $W_2 \in \mathbb{R}^{1 \times H}$, 
$D = 4096$~--- размерность входа, $H = 262144$~--- размерность 
скрытого слоя. Матрица $W_1$ размера $262144 \times 4096$ 
(около 4~ГБ в float32) доминирует в потреблении памяти 
как при прямом проходе, так и при вычислении backward bounds 
в CROWN.

\subsubsection{Реализация Tensor Parallelism}

Для TP-варианта модель была разбита по схеме 
Megatron-LM\cite{MegatronLM2020}:
\begin{itemize}
    \item \textbf{Column Parallel} ($W_1$): весовая матрица 
    разрезана по столбцам (по выходному измерению). Каждый GPU 
    хранит шард $W_1^{(r)} \in \mathbb{R}^{(H/P) \times D}$, 
    где $P$~--- число GPU.
    \item \textbf{Row Parallel} ($W_2$): весовая матрица 
    разрезана по строкам (по входному измерению). Каждый GPU 
    хранит шард $W_2^{(r)} \in \mathbb{R}^{1 \times (H/P)}$, 
    вычисляет частичный результат, после чего выполняется 
    \texttt{AllReduce} (суммирование) для получения полного выхода.
\end{itemize}

Для интеграции с \texttt{auto\_LiRPA}\cite{AutoLiRPARepo} были 
реализованы операторы bound propagation 
(\texttt{BoundLinearTP\_Col}, \texttt{BoundLinearTP\_Row}), 
наследующие от \texttt{BoundLinear} и добавляющие 
\texttt{AllReduce} на этапе обратного распространения границ.

\subsubsection{Методология}

Верификация выполнялась методом CROWN для свойства локальной 
робастности в $L_\infty$-норме с $\varepsilon = 0.01$ на батче 
из $N = 2048$ входных примеров размерности $D = 4096$, 
сгенерированных случайно. Измерялось пиковое потребление 
GPU-памяти (\texttt{torch.cuda.max\_memory\_allocated}).

\subsubsection{Результаты}

\begin{table}[htbp]
\centering
\caption{Пиковое потребление GPU-памяти при верификации методом CROWN 
($D{=}4096$, $H{=}262144$, $N{=}2048$, $\varepsilon{=}0.01$)}
\label{tab:tp-memory}
\begin{tabular}{lccc}
\toprule
Режим & GPU & \texttt{max\_allocated}, МБ & \texttt{max\_reserved}, МБ \\
\midrule
Single GPU & 1$\times$A40 & 26\,826 & 30\,874 \\
TP = 2     & GPU 0        & 13\,513 & 15\,512 \\
TP = 2     & GPU 1        & 13\,513 & 15\,512 \\
\bottomrule
\end{tabular}
\end{table}

Коэффициент снижения памяти составил:
$$k = \frac{26\,826}{13\,513} \approx 1.985,$$
что близко к теоретическому пределу $P = 2$ для двух GPU.
Незначительное отклонение от идеального значения объясняется 
нешардируемыми структурами данных: входной тензор 
$\mathbf{x} \in \mathbb{R}^{N \times D}$, буферы NCCL, а также 
метаданные вычислительного графа, реплицированные на каждом GPU.

\subsection{Эксперимент 2: корректность и звуковость bounds}

\subsubsection{Методология}

Для проверки корректности TP-реализации использовались модели 
MLP разной глубины. Для каждой модели вычислялись bounds двумя 
способами:
\begin{enumerate}
    \item \textbf{Эталон:} стандартный CROWN на одном GPU 
    (полная весовая матрица, без \texttt{AllReduce}).
    \item \textbf{TP=2:} CROWN с Tensor Parallelism на двух GPU 
    (шардированные весовые матрицы с \texttt{AllReduce}).
\end{enumerate}

Проверялись три свойства:
\begin{itemize}
    \item \textbf{Точность:} максимальное отклонение 
    $|lb_{\text{ref}} - lb_{\text{tp}}|$ и 
    $|ub_{\text{ref}} - ub_{\text{tp}}|$.
    \item \textbf{Звуковость:} 
    $lb_{\text{tp}} \leq lb_{\text{ref}}$ и 
    $ub_{\text{tp}} \geq ub_{\text{ref}}$ для всех выходов 
    (TP-bounds должны быть шире эталона).
    \item \textbf{Кросс-ранговая согласованность:} оба GPU 
    должны получить идентичные bounds.
\end{itemize}

Модели PyTorch генерировались программно; модели ONNX были 
загружены из бенчмарка VNN-COMP 2022 
(MNIST-FC)\cite{arXiv2212VNNCOMP2022}. 
Верификация проводилась с $\varepsilon = 0.02$ для 
$L_\infty$-робастности.

\subsubsection{Результаты}

\begin{table}[htbp]
\centering
\caption{Корректность TP=2 CROWN bounds относительно эталона 
(single-GPU CROWN)}
\label{tab:tp-correctness}
\begin{tabular}{lccccc}
\toprule
Модель & Пар Col-Row & $|lb_{\text{diff}}|$ & $|ub_{\text{diff}}|$ 
    & Кросс-ранг & Звуков. \\
\midrule
PyTorch 256$\times$2 & 1 & $2.98 \cdot 10^{-8}$ & $2.98 \cdot 10^{-8}$ & 0 & Да \\
PyTorch 256$\times$4 & 2 & $2.51 \cdot 10^{0}$  & $2.49 \cdot 10^{0}$  & 0 & Да \\
ONNX 256$\times$2    & 1 & $5.96 \cdot 10^{-8}$ & $5.96 \cdot 10^{-8}$ & 0 & Да \\
ONNX 256$\times$4    & 2 & $5.64 \cdot 10^{0}$  & $3.98 \cdot 10^{0}$  & 0 & Да \\
ONNX 256$\times$6    & 3 & $7.92 \cdot 10^{2}$  & $7.50 \cdot 10^{2}$  & 0 & Да \\
\bottomrule
\end{tabular}
\end{table}

\subsubsection{Анализ результатов}

Результаты демонстрируют два ключевых свойства реализации:

\textbf{1. Точность для одной пары Col-Row.} Модели с одной 
парой Column--Row (256$\times$2) дают bounds, совпадающие 
с эталоном с точностью до машинного эпсилона ($\sim 10^{-8}$). 
Это подтверждает арифметическую корректность реализации: 
\texttt{AllReduce} в CROWN backward точно воспроизводит 
результат полного матричного умножения.

\textbf{2. Звуковость с потерей точности для нескольких пар.} 
Модели с несколькими парами (256$\times$4, 256$\times$6) 
дают bounds, которые \textit{шире} эталона (звуковая 
овер-аппроксимация), но менее точны. Потеря точности 
растёт с числом шардированных зон. Это фундаментальное 
следствие архитектуры TP для верификации: промежуточные 
bounds для узлов в шардированных зонах вычисляются методом 
IBP (а не CROWN), поскольку CROWN backward не может 
корректно пересекать границы зон с частичными весовыми 
матрицами. IBP даёт более рыхлые, но гарантированно 
звуковые промежуточные bounds, что приводит к 
овер-аппроксимации в последующих слоях.

\textbf{3. Идеальная кросс-ранговая согласованность.} 
Все модели показывают нулевое расхождение между GPU, 
подтверждая детерминированность вычислений и корректность 
\texttt{AllReduce}-синхронизации.

\subsection{Эксперимент 3: автоматическое шардирование моделей VNN-COMP}

Третий эксперимент проверял работоспособность функции 
автоматического шардирования \texttt{tp\_shard\_bounded\_module} 
на моделях из бенчмарка VNN-COMP 2022\cite{arXiv2212VNNCOMP2022}. 
Модели MNIST-FC (полносвязные сети для классификации MNIST) были 
загружены в формате ONNX, преобразованы в PyTorch с помощью 
\texttt{onnx2pytorch}, обёрнуты в \texttt{BoundedModule} 
и автоматически шардированы без ручного вмешательства.

Функция \texttt{tp\_shard\_bounded\_module} выполняет 
следующие шаги:
\begin{enumerate}
    \item Обход вычислительного графа \texttt{BoundedModule} 
    в топологическом порядке.
    \item Идентификация чередующихся пар \texttt{BoundLinear}.
    \item Замена на \texttt{BoundLinearTP\_Col} / 
    \texttt{BoundLinearTP\_Row} с шардированием весов.
    \item Обновление метаданных графа (размерности выходов, 
    forward-значения) через повторный прямой проход.
    \item Разметка шардированных зон для корректного 
    вычисления промежуточных bounds.
\end{enumerate}

Все три протестированные модели VNN-COMP (256$\times$2, 
256$\times$4, 256$\times$6) были успешно шардированы и прошли 
тест на звуковость (см. Таблицу~\ref{tab:tp-correctness}), 
подтверждая применимость подхода к произвольным 
полносвязным моделям.
